\scanpage{35}
d'où
\[
\Sigma_{12}=X\deleted{-}\setminus(\Sigma_1^+\cup\Sigma_2^+)
\]
\unsure{donc} $\Sigma_{12}=\varnothing$ signifierait que
\[
X=\Sigma_1^+\amalg\Sigma_2^+
\]
et par suite
\[
\unsure{\Pi|\pi}_1=\left\{\Sigma_1^+,\underset{\inserted{\substack{\Vert\\\Sigma_{\unsure{\deleted{2|1}}\unsure{1|2}}^-=\unsure{\deleted{{-}|\Sigma}}X\unsure{\deleted{-}}\setminus\Sigma_1^+}}}{\Sigma_2^+}\right\},
\qquad\Pi_2=\{\Sigma_2^+,\Sigma_1^+\}
\]
alors $\deleted{\unsure{\Pi_2|\Pi_1}}\quad\Pi_1=\Pi_3$ contrairement : l'\unsure{hypothèse|hypothese}.

On trouve donc

\noindent\deleted{\uline{Corollaire}} \inserted{Proposition} \unsure{Soient|Soit} $\Pi_1;\unsure{\deleted{\vert}}\Pi_2$ deux décompositions\\
\makebox[0pt][r]{\rule[-2.7ex]{.4pt}{3.8ex}\hspace{.45em}}binaires de $X$, mutuellement non panachées.\\
Alors il existe un \uline{unique} couple
\[
(\Sigma_1^{\unsure{\top|+}},\Sigma_2^+)\in\Pi_1\times\Pi_2
\]
\noindent\makebox[0pt][r]{\rule[-.35ex]{.4pt}{2ex}\hspace{.45em}}tel que l'on ait
\[
\Sigma_1^+\cap\Sigma_2^+=\varnothing
\]

\begingroup\makeatletter\tagsleft@true\makeatother
Soit alors
\[
C=\Pi_1\amalg\Pi_2=\omega_1\amalg\omega_2
\tag*{(17)}
\]
l'ensemble somme, qui a quatre éléments. On a une application canonique
\[
C\longrightarrow{\unsure{\mathfrak{P}|\mathcal{P}}}_{\unsure{\text{finies}|\text{fermés}|\text{fermé}}}(X)
\tag*{(1\deleted{7}8)}
\]
qui a comme image
\[
\{\Sigma_1^+,\Sigma_1^-,\Sigma_2^+,\Sigma_2^-\},
\tag*{(1\deleted{8}9)}
\]
et laquelle on l'identifie si $(1\deleted{\unsure{7}}8)$ est injective, i.e.
\endgroup
\clearpage
% archive sheet number: 34
